% Introduction to the bishop vs knight chapter, for the 1500-1599 book only.
% There is no shared texts/bishop-vs-knight.tex yet; if one is written, \input it at the end.

\section*{What to look for}
Bishop against knight has little fixed theory, unlike rook endings. What decides these games is a
handful of ideas that come back again and again. Know them, and many positions become easy to judge.

\emph{The knight wins with forks.}\index{Knight fork} The most common tactic in this chapter is a
knight fork, and both sides missed it: the attacker by not looking for it, the defender by walking
into it. Before every king move near a knight, check its squares. A typical trick is to attack a
pawn and fork the pieces that defend it (\xref{Diagram}{mohamed-alagouri-79w} and
\xref{Diagram}{papler-brecelj-61b}).

\emph{The bishop defends best along one diagonal.}\index{One-diagonal principle (bishop)} A bishop
can stop pawns from far away, but only as long as it guards a single diagonal. When it must watch
two pawns on different diagonals, it is torn apart (\xref{Diagram}{daudaravicius-boyajyan-44w}).
The defender, in turn, should arrange things so that one diagonal is enough
(\xref{Diagram}{ciocani-kerimkulov-48w}).

\emph{The knight dislikes the rook's pawn.}\index{Knight against pawns!Rook's pawn} On the edge of
the board the knight has little room. A knight tied to a rook's pawn is a weakness in itself, and
against a passed pawn it often needs a barrier or a blockade to hold
(\xref{Diagram}{papageorgiou-tsolaki-53b} and \xref{Diagram}{pasqualotto-marge-62b}). Against the
bishop, the knight's side often wins with \emph{interference}:\index{Bishop against knight!Interference}
the knight steps onto the bishop's diagonal and cuts it (\xref{Diagram}{luk-davies-66w} and
\xref{Diagram}{allahverdi-mueller-78w}).

\emph{The defender has resources.} Many of these endings can be held even a pawn down. The
defender's tools are a fortress,\index{Fortress} exchanging pawns,\index{Defensive resources!Exchanging pawns}
removing his own weak pawn,\index{Strategy!Remove your own weakness} and in the end giving up the
piece for the last pawns (\xref{Diagram}{solovyev-shilnikov-57b}, \xref{Diagram}{wang-du-55w} and
\xref{Diagram}{ozaltun-vatansever-56w}). When you defend, look for a clear-cut draw first;
when you attack, do not rush.\index{Strategy!Do not rush}
